%=====================================================================
% 12  相互作用量子场论的微扰展开
% 原文：Section 12. Perturbative expansion for interacting QFT（arXiv:2409.03117v1, p. 171--181）
%=====================================================================
\section{相互作用量子场论的微扰展开}

\begin{keytakeaways}
\begin{itemize}
  \item \textbf{物理图像：} 相互作用由拉格朗日量中的高次项（如三次顶点）引入；两点函数的单圈修正给出第一个发散积分，须借依赖于截断的抵消项把发散吸收进质量与耦合，关联函数才可能有有限极限。
  \item \textbf{数学工具：} Feynman 参数化（Feynman 著名公式）、球坐标下的径向积分、表观发散度的幂次计数、按量纲定义的经典标度维数。
  \item \textbf{关键结论：} 按表观发散度的符号把相互作用项分为超可重整、可重整、不可重整三类；可重整理论的重整化只需修改有限多个项，临界维数可直接由顶点的标度维数算出。
  \item \textbf{易错点：} 截断必须以与抵消项量纲相容的方式进入（质量平方、波函数与耦合分别对应不同发散）；表观发散度非负只是发散的充分条件，负值并不保证收敛；二维情形超可重整项有无穷多个。
\end{itemize}
\end{keytakeaways}

\bigskip
\noindent\rule{\textwidth}{0.6pt}
\bigskip

\subsection{量子化的一般策略}
现在我们转入由作用量
$S(\phi):=\int \mathcal L(\phi)dx$ 定义的非自由场论，其中 Minkowski 空间 $V\cong \Bbb R^d$，$\mathcal L(\phi)$ 是局部 Poincaré 不变的拉格朗日量。这类理论量子化的一般策略如下。

{\bf 第 1 步。} 写下该理论的欧氏路径积分关联函数：
$$
\la \phi(x_1)...\phi(x_n)\ra=\int \phi(x_1)...\phi(x_n)e^{-\frac{S_E(\phi)}{\hbar}}D\phi.
$$
用 Feynman 规则计算相应的关于 $\hbar$ 的形式展开（形式幂级数展开）
（正如我们在量子力学 $d=1$ 情形所做的那样）。

{\bf 第 2 步。} 对该形式级数作 Borel 求和，得到对小 $\hbar>0$ 有定义的实际的函数。

{\bf 第 3 步。} 对这些函数作 Wick 转动到 Minkowski 空间，得到 Wightman 关联函数 $W_n$。

{\bf 第 4 步。} 用函数 $W_n$ 定义 Wightman QFT，即提取
Hilbert 空间 $\mathcal H$、Poincaré 群（的双重覆盖）在 $\mathcal H$ 上的表示 $\pi$、
真空向量 $\Omega$ 与场映射 $\phi$。

所有这些步骤都非平凡；虽然第 1 步可以完全严格地完成，但从第 2 步起，严格的实现只对构造性场论中处理过的少数理论已知
（而且对许多拉格朗日量，最终的 Wightman QFT 实际上并不存在）。对大多数物理上有意义的理论，严格完成这些步骤仍是未解问题。本节只讨论第 1 步。

\subsection{$\phi^3$ 理论}

作为贯穿全节的例子，我们使用标量玻色子 $\phi$ 的理论，其欧氏拉格朗日量为
$$
\mathcal L_E(\phi):=\frac{1}{2}((d\phi)^2+m^2\phi^2)+\frac{g}{6}\phi^3,
$$
我们称之为\emph{$\phi^3$ 理论}。该理论是自由标量玻色子理论的形变，
只添加了一个相互作用项 $\frac{g}{6}\phi^3$，
它在 Feynman 演算中对应于一个 3 价顶点。
物理上，该顶点对应于两个粒子碰撞后转化为第三个粒子的相互作用。

我们将取 $\hbar=1$，并考虑关于 $g$ 的幂的形式展开
（由对 $\phi$ 的重新标度，它等价于第 1 步）。

我们来计算自由理论的两点关联函数
$$
\widehat G_0(p)=\frac{1}{p^2+m^2}
$$
在动量表象中的单圈修正。容易看出，该修正由单一 Feynman 图给出

\setlength{\unitlength}{0.435cm}
    \begin{picture}(6,12)(7,-1)

 \put(18,4){\line(1,0){10}}
      \put(18,9){\line(1,0){10}}
      \put(18,4){\line(0,1){5}}
      \put(28,4){\line(0,1){5}}
      %%% make graph
      \put(19,7){\circle*{0.2}}
      \put(19,7){\line(1,0){2.75}}
      \put(21.65,7){\circle*{0.2}}
      \put(22.5,7){\circle{1.75}}
      \put(19,6){\makebox(0,0)[c]{$1$}}
      \put(23.3,7){\line(1,0){2.75}}
      \put(26,6){\makebox(0,0)[c]{$2$}}
      \put(23.35,7){\circle*{0.2}}
      \put(26.05,7){\circle*{0.2}}
      \put(23.5,5){\makebox(0,0)[c]{}}
      \put(27,8){\makebox(0,0)[c]{}}

    \end{picture}

该 Feynman 图的振幅为
$$
A(p)=\frac{g^2}{2(p^2+m^2)^2}\int_V \frac{dq}{(q^2+m^2)((p-q)^2+m^2)}.
$$
若 $d<4$，该积分收敛，可以显式算出。
为此，可使用多元微积分中的下述引理，它在物理学文献中称为\emph{Feynman 著名公式}：

\begin{lemma} 设 $\Delta_n$ 是 $\Bbb R^n$ 中由方程
$$
y_1+...+y_n=1,
$$
定义的 $n-1$ 维单纯形，$dy$ 是 $\Delta_n$ 上体积为 $1$ 的 Lebesgue 测度。则对正数 $a_1,...,a_n$，有
$$
\int_{\Delta_n}\frac{dy}{(a_1y_1+...+a_ny_n)^n}=\frac{1}{a_1...a_n}.
$$
\end{lemma}

\begin{proof} 我们有
$$
\frac{1}{(a_1y_1+...+a_ny_n)^n}=\frac{1}{(n-1)!}\int_0^\infty t^{n-1}e^{-(a_1y_1+...+a_ny_n)t}dt.
$$
于是得到
$$
\int_{\Delta_n}\frac{dy}{(a_1y_1+...+a_ny_n)^n}=\frac{1}{(n-1)!}\int_{\Delta_n}\int_0^\infty t^{n-1}e^{-(a_1y_1+...+a_ny_n)t}dtdy
$$
$$
=\frac{1}{(n-1)!}\int_{t\Delta_n}\int_0^\infty e^{-a_1z_1+...+a_nz_n}dtdz
$$
$$
=\int_{z_1,...,z_n\ge 0}e^{-a_1z_1+...+a_nz_n}dz=\prod_{j=1}^n\int_0^\infty e^{-a_jz_j}dz_j=\frac{1}{a_1...a_n}.
$$
\end{proof}

% [Tutor Note: “Feynman 著名公式”即通常的 Feynman 参数化公式（原文如此命名）。证明中的关键一步是换元 $z_j:=t y_j\ (j=1,\dots,n)$：它把 $(t,y)\in(0,\infty)\times\Delta_n$ 双射到 $\Bbb R_+^n$，且体积元满足 $dz=t^{n-1}dt\,dy$，于是内层积分恰化为 $\int_{\Bbb R_+^n}e^{-\sum_j a_jz_j}dz$。这一恒等式把分母的 $n$ 次幂与单纯形上的积分相联系；实用时常写 $1/(A_1\cdots A_n)=\int_{\Delta_n}dy/(\sum_j y_jA_j)^n$，再把分母凑成 $n=2$ 的平方，从而把圈积分化为对参数 $y$ 的积分。注意有些文献把该式也称为“Feynman 积分”，与本节的“famous formula”是同一件事。]

把 Feynman 著名公式应用于我们的积分，并作变量替换
$q\mapsto q+(1-y)p$，得
$$
\int_V \frac{dq}{(q^2+m^2)((p-q)^2+m^2)}=\int_0^1\int_V \frac{dq}{((1-y)q^2+y(p-q)^2+m^2)^2}dy=
$$
$$
\int_0^1\int_V \frac{dq}{(q^2+M^2(y,p))^2}dy,
$$
其中
$$
M^2(y,p):=y(1-y)p^2+m^2.
$$
现在使用球坐标
$$
\int_V \frac{dq}{(q^2+M^2)^2}=C_d\int_0^\infty \frac{r^{d-1}dr}{(r^2+M^2)^2},
$$
其中 $C_d$ 是 $\Bbb R^d$ 中单位球面的面积。
于是对 $d=2$ 有
$$
\int_V \frac{dq}{(q^2+M^2)^2}=2\pi\int_0^\infty \frac{rdr}{(r^2+M^2)^2}=\pi \int_0^\infty\frac{ds}{(s+M^2)^2}=\frac{\pi}{M^2}.
$$
由此可得
$$
\int_V \frac{dq}{(q^2+m^2)((p-q)^2+m^2)}=\pi\int_0^1 \frac{dy}{y(1-y)p^2+m^2}
$$
$$
=\frac{2\pi}{p^2\sqrt{\frac{4 m^2}{p^2} + 1}}{\rm arccotanh}\sqrt{\frac{4 m^2}{p^2} + 1}.
$$
$d=3$ 的情形可类似计算。

然而，对 $d\ge 4$，我们遇到第一个困难：该积分发散（因为被积函数在无穷远处行为如 $|q|^{-4}$）。更具体地说，对截断 $\Lambda>0$，
定义
$$
A_\Lambda(p):=\frac{g^2}{2(p^2+m^2)^2}\int_{|q|\le \Lambda} \frac{dq}{(q^2+m^2)((p-q)^2+m^2)},
$$
即 $V$ 中半径 $\Lambda$ 的球上的积分。
则当 $d=4$ 时
$$
A_\Lambda(p)\sim \pi^2\frac{g^2}{(p^2+m^2)^2}\log(\tfrac{\Lambda}{m}),\ \Lambda\to \infty
$$
而当 $d>4$ 时
$$
A_\Lambda(p)\sim C_d\frac{g^2}{2(d-4)(p^2+m^2)^2}\Lambda^{d-4},\ \Lambda\to \infty
$$
补救这一困难的办法是在拉格朗日量中添加一个依赖于 $\Lambda$ 的项，称为\emph{抵消项}；它在 $\Lambda\to\infty$ 时发散，但会抵消这个发散，其意义是：当在球 $|q|\le \Lambda$ 上积分时，该积分在 $\Lambda\to\infty$ 时有有限极限。

例如，取 $d=4$。此时模去 $g^3$，在截断 $\Lambda$ 下算出的动量表象两点函数形如
$$
\widehat G_{\Lambda,m^2}(p)=\frac{1}{p^2+m^2}+\pi^2\frac{g^2}{(p^2+m^2)^2}\log(\tfrac{\Lambda}{m})+...
$$
（这里我们明确写出 $\widehat G$ 对 $m^2$ 的依赖，因为我们即将改变它）。
让我们尝试用 $m^2+Kg^2\log(\tfrac{\Lambda}{m})$（其中 $K$ 为常数）替换参数 $m^2$ 来修正发散。于是有
$$
\widehat G_{\Lambda,m^2+Kg^2\log(\tfrac{\Lambda}{m})}(p)=\frac{1}{p^2+m^2+Kg^2\log(\tfrac{\Lambda}{m})}+\pi^2\frac{g^2}{(p^2+m^2)^2}\log(\tfrac{\Lambda}{m})+...
$$
$$
=\frac{1}{p^2+m^2}+(\pi^2-K)\frac{g^2}{(p^2+m^2)^2}\log(\tfrac{\Lambda}{m})+...
$$
其中忽略高于 $g^2$ 阶的项。因此要抵消发散，应取
$K=\pi^2$，即把拉格朗日量换成
$$
\mathcal L_{E,\Lambda}:=\frac{1}{2}((d\phi)^2+(m^2+\pi^2g^2\log(\tfrac{\Lambda}{m}))\phi^2)+\frac{g}{6}\phi^3.
$$
对这个拉格朗日量，若以截断 $\Lambda$ 积分，则两点函数模去 $g^2$ 后当
$\Lambda\to \infty$ 有有限极限，由下式给出
$$
\widehat G(p)=\frac{1}{p^2+m^2}+\frac{g^2}{2(p^2+m^2)^2}I(p),
$$
其中
$$
I(p)=\lim_{\Lambda\to \infty} \left(\int_{\Bbb R^4} \frac{dq}{(q^2+m^2)((p-q)^2+m^2)}-2\pi^2\log(\tfrac{\Lambda}{m})\right).
$$
利用 Feynman 著名公式很容易算出这个极限。即，类似 $d<4$ 情形计算可得
$$
I(p)=\int_0^1 I(p,y)dy,\ I(p,y):=\lim_{\Lambda\to \infty} \left(\int_0^\Lambda \frac{r^3dr}{(r^2+M^2(y,p))^2}-2\pi^2\log (\tfrac{\Lambda}{m})\right).
$$
于是
$$
I(p,y)=2\pi^2(\log m-\frac{1}{2}(1+\log(y(1-y)p^2+m^2))).
$$
因此
$$
I(p)=2\pi^2\left(\frac{1}{2}+\sqrt{\frac{4 m^2}{p^2} + 1}\cdot {\rm arccotanh}\sqrt{\frac{4 m^2}{p^2} + 1}\right).
$$

对 $d>4$，计算变得更复杂。
即，虽然 $d=5$ 时
$$
A_\Lambda(p)\sim C_5\frac{g^2}{2(p^2+m^2)^2}\Lambda+O(1), \lambda\to \infty,
$$
故程序相同，只是质量参数需改为
$m^2\mapsto m^2+K\Lambda$，但到 $d=6$ 时我们必须对发散积分作更精细的展开：
$$
A_\Lambda(p)\sim \frac{g^2}{4(p^2+m^2)^2}(C_6\Lambda^2+Cp^2\log(\tfrac{\Lambda}{m})+O(1)),\ \Lambda\to \infty
$$
我们可以通过质量修正
$m^2\mapsto m^2+K\Lambda^2$ 抵消最奇异项 $C_6\Lambda^2$，但此后仍有对数发散
$Cp^2\log(\tfrac{\Lambda}{m})$，它依赖于 $p$。为消除这个发散，必须用抵消项修改
拉格朗日量中 $\frac{1}{2}(d\phi)^2$ 的系数，把它由 $1$
改为 $1+C'g^2\log(\tfrac{\Lambda}{m})$，其中 $C'$ 是适当的常数。此外我们发现，三点函数的单圈修正
是对数发散的：
相应的动量表象贡献（至多差一个标度）为
$$
\frac{g^3}{\prod_{j=1}^3(p_j^2+m^2)}J(p_1,p_2,p_3)\delta(p_1+p_2+p_3),
$$
其中
$$
J(p_1,p_2,p_3)=\int_V \frac{dq}{(q^2+m^2)((q-p_1)^2+m^2)((q-p_1-p_2)^2+m^2)}
$$
（对 $p_1+p_2+p_3=0$），该积分发散，在半径
$\Lambda$ 的球上计算时行为如 $\log \Lambda$。因此为消除这个发散，必须用抵消项修改拉格朗日量中 $\frac{1}{6}\phi^3$
的系数，把它由 $g$ 改为 $g+C''g^3\log(\tfrac{\Lambda}{m})$。

我们开始看到\emph{重整化理论}的主要思想，它使我们能够对所有微扰阶的 Feynman 图所产生的发散积分作正规化。这一思想是：拉格朗日量中的系数实际上\textbf{并非}有意义的物理量——它们只是依赖于计算所用标度（截断）$\Lambda$ 的数学参数，并可能在 $\Lambda\to \infty$（称为\emph{紫外极限}，因为 $\Lambda$ 具有振动频率的意义）时发散。真正有意义的是答案，即关联函数 $\la \phi(x_1)...\phi(x_n)\ra$（或其在动量表象中的 Fourier 变换）。这个答案依赖于某些参数，它们才是理论的真正参数。因此，拉格朗日量中的系数必须作相应调整，使
答案在 $\Lambda\to \infty$ 时有有限极限。我们最终得到的具体答案将依赖于调整程序，但在好的情形（称为可重整）下会落在一个漂亮的普遍族中（通常依赖有限多个参数，但并非总是如此）。

\subsection{超可重整、可重整与不可重整理论}
让我们更系统地讨论这一点。考虑具有一般拉格朗日量的标量玻色子理论。给定
Feynman 图 $\Gamma$，有相应的动量表象中的 Feynman 积分 $I_\Gamma$，它是实向量空间上有理体积形式的积分。我们可以定义\emph{表观发散度} $D(\Gamma)$ 为该形式的分子次数（其中坐标微分的次数记为 $1$）减去分母次数。显然，若 $D(\Gamma)\ge 0$ 则积分发散。注意逆命题不成立：若 $d(\Gamma)<0$，积分仍可能发散。

我们来计算 $D(\Gamma)$。分母的次数容易计算：
它恰为 $2e(\Gamma)$，其中 $e(\Gamma)$ 是 $\Gamma$ 的内边数
（因为每条边贡献一个传播子，即二次函数的倒数）。
另一方面，对 $V$ 的积分次数等于圈数，即
$d(e(\Gamma)-v(\Gamma)+1)$，其中 $v(\Gamma)$ 是内顶点数。
最后，拉格朗日量中含 $\phi$ 导数的项把相应导数个数贡献给分子次数。由此可得
$$
D(\Gamma)=(d-2)e(\Gamma)-dv(\Gamma)+d+N,
$$
其中 $N$ 是各顶点单项式中导数总数。特别地，当没有导数时，有
$$
D(\Gamma)=(d-2)e(\Gamma)-dv(\Gamma)+d.
$$
这表明 $D(\Gamma)$ 可按各顶点贡献之和计算：把位于全内顶点（所有边都是内边的顶点）的微分单项式 $\Phi$ 的次数 $D(\Phi)$ 定义为该顶点对 $D(\Gamma)$ 的贡献。
事实上，每个 $\Phi$ 贡献
$$
D(\Phi)=\frac{d-2}{2}e(\Phi)-d+N_\Phi,
$$
其中 $e(\Phi)$ 是 $\Phi$ 的边数（即它关于 $\phi$ 的次数），
而 $N_\Phi$ 是 $\Phi$ 中导数的个数。

我们看到更自然的量是
$$
[\Phi]:=D(\Phi)+d,
$$
因为它具有乘性：
$$
[\Phi_1\Phi_2]=[\Phi_1][\Phi_2].
$$
这并不奇怪，因为 $\Phi$ 带有体积因子 $dx$，
故 $D(\Phi)$ 实际上是 $\Phi dx$ 的标度维数；因此
要得到 $\Phi$ 的标度维数，需加上
$d$（因为 $dx$ 的标度维数是 $-d$）。
这引出了

\begin{definition} 数 $[\Phi]$ 称为微分单项式 $\Phi$ 的\emph{经典标度维数}。
\end{definition}

于是对 Feynman 图 $\Gamma$ 有

\begin{equation}\label{degfo} 
D(\Gamma)=d-\frac{k(d-2)}{2}+\sum_\Phi D(\Phi),
\end{equation}
其中 $k$ 是 $\Gamma$ 的外顶点数。

例如，对 $\Phi=\phi^n$ 得到
$$
D(\phi^n)=\tfrac{n}{2}(d-2)-d=(\tfrac{n}{2}-1)d-n,
$$
每个导数使次数增加 $1$，因此例如
$$
D(\phi^{n-2}(d\phi)^2)=(\tfrac{n}{2}-1)(d-2).
$$
于是对 $k$ 点函数的单圈 Feynman 图 $\Gamma$（带 $k$ 条leg的圈），有
$$
D(\Gamma)=d-\frac{k}{2}(d-2)+kD(\phi^3)=d-\frac{k}{2}(d-2)+\frac{k}{2}(d-6)=d-2k. 
$$

\begin{definition} 设 $\Phi$ 是 $\phi$ 的微分单项式。若
$D(\Phi)<0$，则称 $\Phi$ 为\emph{超可重整}；若 $D(\Phi)=0$，则称\emph{可重整}（或临界）；若
$D(\Phi)>0$，则称\emph{不可重整}。
\end{definition}

因此超可重整项改善收敛性，可重整项不影响收敛性，而
不可重整项使其变坏。

\begin{example} 1. 动能项 $(d\phi)^2$ 有 $D=0$，故可重整；
事实上，按定义这在任何 QFT 中都如此。
注意这可用来容易地
计算单项式的经典标度维数。即
$[(d\phi)^2]=D((d\phi)^2)+d=d$，故
$2[\phi]+2=d$，即 $[\phi]=\frac{d-2}{2}$。
利用乘性，现在即可立即算出任意 $\Phi$ 的 $[\Phi]$。

2. 质量项 $\phi^2$ 有 $D=-2$，故超可重整。
项 $\phi^3$ 有 $D=\frac{1}{2}d-3$，故当 $d<6$ 时超可重整，
$d=6$ 时可重整，$d>6$ 时不可重整。
\end{example}

\begin{definition} 称拉格朗日量为

$\bullet$ 超可重整的，若其除动能项外的所有项都超可重整；

$\bullet$ 可重整的（或临界的），若其所有项至多可重整，且至少有一个可重整的非动能（即相互作用）项；

$\bullet$ 不可重整的，若它含有不可重整项。
\end{definition}

显然，每个拉格朗日量恰属于这三种类型之一。

\begin{proposition} 
(i) 若拉格朗日量超可重整，则相应 Feynman 图的表观发散度有上界，且对任意给定的外线数只有有限多个表观发散图；此外，若 $d>2$，则表观发散图总共只有有限多个。

(ii) 若拉格朗日量可重整，则固定外线数的表观发散图有无穷多个，但这些图的表观发散度仍有上界。

(iii) 若拉格朗日量不可重整，则固定外线数的图的表观发散度无上界。
\end{proposition}

\begin{proof} 这由公式 \eqref{degfo} 即知。
\end{proof}

这意味着对不可重整的拉格朗日量，发散积分的正规化必定失控。即：若想正规化
表观发散度无上界的图，就必须引入含任意多导数的抵消项，重整化后的拉格朗日量将不再依赖 $\phi$ 的有限多个导数。

另一方面，若拉格朗日量可重整，则当 $d>2$ 时，重整化程序中只需修改有限多个
项；它们正是拉格朗日量中可能的超可重整项与可重整项。该程序对所有微扰阶都有效，这是一个相当非平凡的事实，我们在此不予证明；但其结果是
一个有限参数族的微扰 QFT。

在二维情形还有一个额外特征——
拉格朗日量中有无穷多个（超）可重整项；但它们至多含两个导数。

最后，在超可重整情形，重整化程序在有限步内完成。

\subsection{若干重要 QFT 的临界维数}\label{critthe}

对由拉格朗日量定义的相互作用 QFT，理论只在小维数下（超）可重整，
当维数增大时变为不可重整。若某理论在某个维数 $d$ 可重整、
而在更大维数不可重整，则称 $d$ 为该理论的\emph{临界维数}。

\subsubsection{标量玻色子}
例如，由于 $D(\phi^n)=(\tfrac{n}{2}-1)d-n$，对标量玻色子，
项 $\phi^n$ 是（超）可重整的当且仅当
$d\le \frac{2n}{n-2}$。于是在（超）可重整理论中，
项 $\phi^3$ 只能出现在 $d\le 6$ 时，$\phi^4$
只能出现在 $d\le 4$ 时，$\phi^5$ 与 $\phi^6$ 只能出现在 $d\le 3$ 时。
又因 $D(\phi^{n-2}(d\phi)^2)=(\tfrac{n}{2}-1)(d-2)$，这类 $n>2$ 的项除非 $d=2$
不能出现在（超）可重整理论中。导数更多时情况更糟。
于是得到

\begin{proposition} 对标量玻色场 $\phi$，最一般的（超）可重整
非二次 Poincaré 不变拉格朗日量为（至多差一个整体标度）：

$\bullet$ $d>6$：无；

$\bullet$: $d=5,6$：$\mathcal L=\frac{1}{2}(d \phi)^2+P_3(\phi)$；

$\bullet$: $d=4$：$\mathcal L=\frac{1}{2}(d \phi)^2+P_4(\phi)$；

$\bullet$: $d=3$：$\mathcal L=\frac{1}{2}(d \phi)^2+P_6(\phi)$；

$\bullet$: $d=2$：$\mathcal L=\frac{1}{2}g(\phi)(d \phi)^2+U(\phi)$，

其中 $P_m$ 是次数为 $m$ 的多项式，$U$ 与 $g$ 是任意
（实解析）函数。
\end{proposition}

注意，不妨假设 $P_m$ 不含常数项与一次项。于是拉格朗日量为 $\frac{1}{2}(d \phi)^2+P_m(\phi)$ 的理论的参数个数为 $m-1$（即 $P_m$ 的系数）。

\subsubsection{费米子}

回忆对费米场 $\psi$，动能项
形如 $(\psi,\bold D\psi)$。这意味着
$$
2[\psi]+1=d,
$$
亦即
$$
[\psi]=\frac{d-1}{2},
$$
它总是正的。因此对质量项 $(\psi,M\psi)$ 有 $D=-1$，它们超可重整。
在二次项之外，我们看到对 $d\ge 2$，$\psi$ 中唯一可能（超）可重整的项
是一般形状为 $\psi^{2k}$ 的项，且
$$
D(\psi^{2k})=2k[\psi]-d=k(d-1)-d=(k-1)(d-1)-1.
$$
它（超）可重整的唯一情形是 $d=2$ 且 $k=2$，即
项 $\psi^4$，此时 $D=0$（临界）。这类项确实出现在所谓的\emph{Gross-Neveu 模型}中。

对 $d>2$，可重整拉格朗日量中的任何费米项因此必须是费米子的二次项。
但它可以含有其他（玻色）场作为因子。例如，
$[\phi^n\psi^2]=n\frac{d-2}{2}+d-1$，故
$$
D(\phi^n\psi^2)=n\frac{d-2}{2}-1.
$$
这表明在 3 维可以有项
$\phi\psi^2$（Yukawa 相互作用）与 $\phi^2\psi^2$，而在 $4$ 维只能有
Yukawa 项 $\phi\psi^2$，当 $d>4$ 时不存在可能的（超）可重整项。

\subsubsection{规范场论}
当 $\phi$ 取向量值（即具有任意多个分量）时，类似结果成立。这使我们能处理另一个重要例子，即\emph{规范场论}。

回忆第 \ref{manif} 小节：为定义规范场论，固定紧 Lie 群 $G$（例如 $U(n)$），场是 $V$ 上主 $G$-丛 $P$ 上的联络 $\nabla$。由于这类丛都平凡，可把 $\nabla$ 视为取值于 $\g={\rm Lie}G$ 的 1-形式 $A$；即 $\nabla_A=d+A$。$\nabla_A$ 的曲率由公式
$$
F_A=dA+\frac{1}{2}[A,A],
$$
给出，纯规范场论的拉格朗日量为
$$
\mathcal L:=\int_V |F_A|^2dx.
$$
如第 \ref{manif} 小节所述，这里的微妙之处在于 $A$ 只在规范变换
$\nabla_A\mapsto g^{-1}\nabla_Ag$，即
$A\mapsto g^{-1}dg+g^{-1}Ag$ 下考虑，其中 $g: V\to G$
是在无穷远处有指定行为的光滑函数，
但这对临界维数的讨论无关紧要。

若 $G$ 交换（例如 $G=U(1)$），则拉格朗日量是二次的，该理论是自由的（即
无物质时的量子电动力学，即 Maxwell 方程的量子化）。
该理论在所有维数下满足 Wightman 公理，其 Wightman 函数
可类似标量玻色子的情形显式算出。

然而，若 $G$ 非交换（例如标准模型中弱相互作用对应 $G=SU(2)$、强相互作用对应 $G=SU(3)$），则拉格朗日量不是二次的，运动方程不是线性的（它们是\emph{Yang-Mills 方程}）。把 $A$ 视为（向量值的）玻色子，
我们看到拉格朗日量中的非二次项
具有示意形式 $A^2dA$ 与 $A^4$。这些项的次数为
$\frac{1}{2}(d-4)$ 与 $d-4$，故该理论在 $4$ 维是临界的
（物理情形！），在低维超可重整，但对 $d>4$ 不可重整。
注意我们有向量玻色子而非
一族标量玻色子（在 $\bold P$ 的作用下）这一事实对维数计数无关紧要。

还要注意，在 $d\le 4$ 维我们也可考虑
含项 $(\nabla_A\phi)^2$ 或 $(\psi,\bold D_A\psi)$ 的可重整拉格朗日量，其中 $\phi$
是标量、$\psi$ 是取值于相伴丛 $P\times_G \rho$ 的旋量，$\rho$ 是 $G$ 的有限维
表示（容易验证出现的所有单项式都有 $D\le 0$）。这类项确实出现在
标准模型中；最简单的例子是 $(\psi,\bold D_A\psi)$，其中
$A$ 是 $U(1)$-联络、$\psi$ 是取值于 $U(1)$ 的定义表示（tautological/defining representation，即 $U(1)$ 按相位乘法作用的一维表示）的旋量，对应于电子。

\subsubsection{$\sigma$ 模型}

\emph{$\sigma$ 模型}是取值于 Riemann 流形 $M$ 的标量玻色子理论。于是场是映射 $\phi: V\to M$，拉格朗日量为
$\mathcal L=\frac{1}{2}(d\phi)^2$，它在局部坐标下形如
$$
\mathcal L=\frac{1}{2}\sum_{i,j=1}^{\dim M}g_{ij}(\phi)d\phi^id\phi^j,
$$
其中 $g_{ij}$ 是 $M$ 上的 Riemann 度量。
我们也可以添加势 $U(\phi)$，其中 $U$ 是 $M$ 上的光滑函数。
由上述计算，对非常数度量，该拉格朗日量
只在维数 $d=2$ 时可重整，但此时 $g_{ij}$ 与 $U$ 可任意。

\subsubsection{引力}

引力理论（广义相对论）是取值于对称张量 $S^2V^*$ 的玻色场 $h(x)$
的理论；即 $V$ 上的 Minkowski 度量经
$g=g_0+h$ 微扰，其中 $g_0$ 是标准
Minkowski 度量。广义相对论的拉格朗日量为
$$
\mathcal L=R(g)
$$
其中 $R$ 是度量 $g$ 的标量曲率。由于曲率用度量的二阶导数表示，在相差一个标度的意义下，它可以示意地写成 $h$ 的形式
$$
\mathcal L=(dh)^2+...
$$
其中省略号代表至多含 $h$ 的两个导数的项。因此该拉格朗日量的一般形状（就计算经典标度维数而言）与 $\sigma$ 模型相同；故该理论只在二维可重整。这是至今未能把引力纳入标准模型的主要原因之一，而标准模型存在于 $4$ 维时空。

\begin{remark} 
我们在第 \ref{compop} 小节已经看到：
即使在自由量子场论中，像 $\phi^2(x)$ 这样的复合算子也不自动有定义，
需要正规序程序来正规化它们。在相互作用 QFT 中更是如此。

结果表明，在临界微扰 QFT 中，正规序程序、复合算子与算子乘积展开
可以用重整化理论类似自由情形地定义。我们在此不讨论，读者可参阅 \cite{QFS} 第 1 卷 p. 452。
\end{remark}

% [Tutor Note: 幂次计数与截断——计算 $D(\Gamma)=(d-2)e-dv+d+N$ 时，圈数 $L=e-v+1$ 给出积分次数 $dL$；每条内边（传播子）贡献分母次数 $2$；每个内顶点处的导数给出分子次数。对 $e=v$ 的图（如 $\phi^3$ 两点泡泡图与三点三角图），$d=4$ 时 $D$ 化为 $4-2v+N$，故两点函数仅需质量抵消项；$d=6$ 时化为 $6-2e+N$，出现 $p^2\log\Lambda$ 型波函数重整化。核对表的独立办法：把 $k$ 点函数的一圈图（带 $k$ 条leg的圈）代入公式 \eqref{degfo} 应得 $D=d-2k$（原文已作此检验）。又，截断必须以与量纲相容的方式进入抵消项：$d=4$ 时质量平方的修正为 $g^2\log(\Lambda/m)$ 乘以常数、波函数为 $g^2\log(\Lambda/m)$、三点耦合为 $g^3\log(\Lambda/m)$，分别对应 $D\ge0$ 的三个顶点类型，不能混用。]

% [Tutor Note: 原文疑误（均照录）——(1) $d=5$ 的渐近式中原文写作 “$\lambda\to \infty$”（小写 lambda），应为 $\Lambda\to\infty$；(2) 12.3 节 “if $d(\Gamma)<0$, the integral may still diverge” 中的 $d(\Gamma)$ 应为 $D(\Gamma)$（表观发散度）；(3) 费米子段 “super-renomalizable” 与规范段 “he subtlety here” 各有一处拼写遗漏；(4) 12.2 节把 Feynman 参数化公式称作 “Feynman famous formula”，非标准名称，译文保留“著名公式”并批注。]

% [Tutor Note: 经典标度维数——$[\Phi]:=D(\Phi)+d$ 的乘性来自 $\Phi dx$ 的标度维数：$D(\Phi)$ 是 $\Phi dx$ 的标度维数，而 $[dx]=-d$，故 $[\Phi]=D(\Phi)+d$。由动能项可定出 $[\phi]=\frac{d-2}{2}$，于是 $[\phi^n]=n\frac{d-2}{2}$、$D(\phi^n)=(\frac{n}{2}-1)d-n$。临界维数即 $D(\Phi)=0$ 的维数：$\phi^3$ 为 $d=6$、$\phi^4$ 为 $d=4$、$\phi^5,\phi^6$ 为 $d=3$；含两个导数的项 $\phi^{n-2}(d\phi)^2$ 的 $D=(\frac{n}{2}-1)(d-2)$，故除 $d=2$ 外都不可重整。$d=2$ 时拉格朗日量中允许无穷多个（超）可重整项（因 $D$ 随 $n$ 降低），这正是二维场论（尤其 $\sigma$ 模型）参数无穷多的来源，与第 13 章共形场论直接相关。]

% [Tutor Note: 规范场与引力——$A^2dA$ 的 $D=\frac{d-4}{2}$、$A^4$ 的 $D=d-4$（用 $[A]=\frac{d-2}{2}$、$e=3,N=1$ 与 $e=4,N=0$ 代入 $D(\Phi)=\frac{d-2}{2}e-d+N$），故 Yang-Mills 临界维数为 $4$，与实验维数吻合；这也是规范场论远比 $\phi^4$ 受青睐的结构性原因。引力写成 $g=g_0+h$ 后 $\mathcal L=(dh)^2+\cdots$，与 $\sigma$ 模型同型，故只在 $d=2$ 可重整（$d=4$ 时引力按幂次计数不可重整），这正是把引力纳入标准模型的主要障碍。注意原文强调“向量玻色子 vs 标量多重态”对维数计数无影响，因为度规通过 $[A]$ 或 $[\phi]$ 只进入导数个数。]

% [Tutor Note: 本章小结——微扰展开给出关联函数的形式幂级数；发散表现为圈积分在 $|q|\to\infty$ 的幂次发散，抵消项以 $\Lambda$ 依赖的系数出现，其形式由表观发散度决定：$d=4$ 的 $\phi^3$ 理论只需质量、波函数与耦合三类抵消项。超可重整/可重整/不可重整的分野由顶点的经典标度维数 $[\Phi]$ 与 $D(\Phi)$ 的符号给出；可重整性使得抵消项只涉及拉格朗日量中原有的有限多个项。本章只完成第 1 步（Feynman 规则），Borel 求和与 Wick 转动（第 2、3 步）以及 Wightman 重建（第 4 步）所需的严格解析工具未在此展开；原文建议参阅 \cite{QFS}（尤其第 1 卷第 452 页关于复合算子与 OPE 的讨论）。]

